% ── Answer Key: Practice 1 — Set Theory ──────────────────────────
% This file is meant to be \input{} into a master document.
% It assumes amsmath, amssymb, amsthm are loaded and
% \newcommand{\comp}[1]{\overline{#1}}, \newcommand{\symdiff}{\mathbin{\triangle}} exist.

\section*{Practice 1: Set Theory}

%======================================================================
\subsection*{Problem 1}
What does $\{\varnothing\}$ represent and how many elements does it have?

$\{\varnothing\}$ is the set whose only element is the empty set. It has $\boxed{1}$ element.

%======================================================================
\subsection*{Problem 2}
Which statements are true?

\begin{itemize}
  \item[2.1.] $\varnothing \in \{\varnothing, \{\varnothing\}\}$ --- \textbf{TRUE}. The set lists $\varnothing$ as an element.
  \item[2.2.] $\{\varnothing\} \in \{\{\varnothing\}\}$ --- \textbf{TRUE}. The sole element of $\{\{\varnothing\}\}$ is $\{\varnothing\}$.
  \item[2.3.] $\varnothing \in \{\{\varnothing\}\}$ --- \textbf{FALSE}. The sole element is $\{\varnothing\} \neq \varnothing$.
\end{itemize}

%======================================================================
\subsection*{Problem 3}
How many elements does $\{\{\varnothing\}\}$ have?

$\boxed{1}$. Its only element is $\{\varnothing\}$.

%======================================================================
\subsection*{Problem 4}
How many elements do the following sets have?

\begin{itemize}
  \item[4.1.] $\{1,2,1\}$: $\boxed{2}$ (duplicates are ignored; the set is $\{1,2\}$).
  \item[4.2.] $\{1,2,\{1,2\}\}$: $\boxed{3}$ (elements: $1$, $2$, $\{1,2\}$).
  \item[4.3.] $\{\{2\}\}$: $\boxed{1}$ (element: $\{2\}$).
  \item[4.4.] $\{1,\{1\}\}$: $\boxed{2}$ (elements: $1$ and $\{1\}$; these are distinct).
  \item[4.5.] $\{1,\varnothing\}$: $\boxed{2}$ (elements: $1$ and $\varnothing$).
  \item[4.6.] $\{\{\varnothing\},\varnothing\}$: $\boxed{2}$ (elements: $\{\varnothing\}$ and $\varnothing$).
  \item[4.7.] $\{\{\varnothing\},\{\varnothing\}\}$: $\boxed{1}$ (duplicate; the set is $\{\{\varnothing\}\}$).
  \item[4.8.] $\{x \mid 2x=1\}$ where $x \in \mathbb{R}$: $\boxed{1}$ (the set is $\{\tfrac{1}{2}\}$).
\end{itemize}

%======================================================================
\subsection*{Problem 5}
Is $\{1,2\} \in \{1,2,\{1,3\},\{1,2,3\}\}$?

\textbf{NO.} The elements of the right-hand set are $1$, $2$, $\{1,3\}$, and $\{1,2,3\}$. None of these equals $\{1,2\}$.

%======================================================================
\subsection*{Problem 6}
Give examples of sets that are elements of themselves.

In standard ZFC set theory (with the Axiom of Regularity), no set is an element of itself. However, in naive set theory or theories without regularity, the ``set of all sets'' $V = \{x \mid x = x\}$ would satisfy $V \in V$. Another classic example: if we define $S$ as ``the set of all sets that are not sets of numbers,'' then $S \in S$ (Russell-type construction). In ZFC, \textbf{no such set exists}.

%======================================================================
\subsection*{Problem 7}
Give examples of $A, B, C$ with $A \in B$, $B \in C$, but $A \notin C$.

Let $A = 1$, $B = \{1,2\}$, $C = \{\{1,2\},3\}$.
Then $1 \in \{1,2\}$, $\{1,2\} \in \{\{1,2\},3\}$, but $1 \notin \{\{1,2\},3\}$ since the elements of $C$ are $\{1,2\}$ and $3$.

%======================================================================
\subsection*{Problem 8}
Prove: $\{\{a\},\{a,b\}\} = \{\{c\},\{c,d\}\} \iff a=c$ and $b=d$.

\textit{Proof.}
($\Leftarrow$) Trivial by substitution.
($\Rightarrow$) From the equality, $\{a\}$ must equal either $\{c\}$ or $\{c,d\}$.
If $\{a\} = \{c,d\}$ then $c = d = a$, and then $\{a,b\} = \{c\} = \{a\}$ forces $b = a = c = d$.
If $\{a\} = \{c\}$ then $a = c$, and the remaining element gives $\{a,b\} = \{c,d\} = \{a,d\}$, so $b = d$. \qed

%======================================================================
\subsection*{Problem 9}
Prove three set equalities/inclusions.

\textbf{(a)} $\{x \in \mathbb{Z} \mid \exists y,\, x = 6y\} = \{x \in \mathbb{Z} \mid \exists u,v,\, x = 2u \text{ and } x = 3v\}$.

\textit{Proof.} ($\subseteq$) If $x = 6y$ then $x = 2(3y)$ and $x = 3(2y)$.
($\supseteq$) If $x = 2u$ and $x = 3v$, then $x$ is divisible by both 2 and 3, hence by $\mathrm{lcm}(2,3) = 6$. \qed

\textbf{(b)} $\{x \in \mathbb{R} \mid \exists y,\, x = y^2\} = \{x \in \mathbb{R} \mid x \ge 0\}$.

\textit{Proof.} ($\subseteq$) $x = y^2 \ge 0$.
($\supseteq$) If $x \ge 0$, take $y = \sqrt{x}$, then $x = y^2$. \qed

\textbf{(c)} $\{x \in \mathbb{Z} \mid \exists y,\, x = 6y\} \subset \{x \in \mathbb{Z} \mid \exists y,\, x = 2y\}$.

\textit{Proof.} If $x = 6y$ then $x = 2(3y)$, so $\subseteq$ holds. The inclusion is strict since $2 = 2 \cdot 1$ is in the right set but $2$ is not a multiple of 6. \qed

%======================================================================
\subsection*{Problem 10}
Prove transitivity properties of $\subseteq$ and $\subset$.

\textbf{(a)} If $A \subseteq B$ and $B \subseteq C$, then $A \subseteq C$.

\textit{Proof.} Let $x \in A$. Then $x \in B$ (since $A \subseteq B$), then $x \in C$ (since $B \subseteq C$). \qed

\textbf{(b)} If $A \subseteq B$ and $B \subset C$, then $A \subset C$.

\textit{Proof.} By (a), $A \subseteq C$. Since $B \subset C$, there exists $z \in C \setminus B$. If $A = C$, then $z \in A \subseteq B$, contradiction. So $A \neq C$, giving $A \subset C$. \qed

\textbf{(c)} If $A \subset B$ and $B \subseteq C$, then $A \subset C$.

\textit{Proof.} By (a), $A \subseteq C$. Since $A \subset B$, there exists $z \in B \setminus A$, and $z \in C$ since $B \subseteq C$. So $z \in C \setminus A$, hence $A \neq C$. \qed

\textbf{(d)} If $A \subset B$ and $B \subset C$, then $A \subset C$.

\textit{Proof.} Follows from (c) since $B \subset C$ implies $B \subseteq C$. \qed

%======================================================================
\subsection*{Problem 11}
Give examples of $A, B, C, D, E$ with $A \subseteq B$, $B \in C$, $C \subseteq D$, $D \subseteq E$.

$A = \{1\}$, $B = \{1,2\}$, $C = \{\{1,2\},3\}$, $D = \{\{1,2\},3,4\}$, $E = \{\{1,2\},3,4,5\}$.

%======================================================================
\subsection*{Problem 12}
Which statements are true for all sets $A, B, C$?

\begin{itemize}
  \item[(a)] If $A \notin B$ and $B \notin C$, then $A \notin C$. --- \textbf{FALSE.}
    Counterexample: $A = 1$, $B = \{2\}$, $C = \{1\}$. Then $1 \notin \{2\}$ and $\{2\} \notin \{1\}$, but $1 \in \{1\}$.
  \item[(b)] If $A \neq B$ and $B \neq C$, then $A \neq C$. --- \textbf{FALSE.}
    Counterexample: $A = C = \{1\}$, $B = \{2\}$.
  \item[(c)] If $A \in B$ and $B \subseteq C$, then $A \in C$. --- \textbf{TRUE.}
    Since $A \in B$ and every element of $B$ is in $C$, we get $A \in C$.
  \item[(d)] If $A \subseteq B$ and $B \subseteq C$, then $C \subseteq A$? --- \textbf{FALSE.}
    Counterexample: $A = \varnothing$, $B = C = \{1\}$.
  \item[(e)] If $A \subseteq B$ and $B \in C$, does $A \notin C$? --- \textbf{FALSE} (not necessarily).
    Counterexample: $A = \{1\}$, $B = \{1,2\}$, $C = \{\{1\},\{1,2\}\}$. Then $A \subseteq B$, $B \in C$, and also $A \in C$.
\end{itemize}

%======================================================================
\subsection*{Problem 13}
Give examples of sets where every element is a subset of the set (transitive sets).

$\varnothing$ (vacuously true). Also $\{\varnothing\}$: its only element $\varnothing \subseteq \{\varnothing\}$.
Also $\{\varnothing, \{\varnothing\}\}$: $\varnothing \subseteq \{\varnothing, \{\varnothing\}\}$ and $\{\varnothing\} \subseteq \{\varnothing, \{\varnothing\}\}$.
In general, any ordinal number is a transitive set.

%======================================================================
\subsection*{Problem 14}
List the elements of $2^A$.

\begin{itemize}
  \item[14.1.] $A = \varnothing$: $\boxed{2^A = \{\varnothing\}}$ (1 element).
  \item[14.2.] $A = \{x\}$: $\boxed{2^A = \{\varnothing, \{x\}\}}$ (2 elements).
  \item[14.3.] $A = \{1,a\}$: $\boxed{2^A = \{\varnothing, \{1\}, \{a\}, \{1,a\}\}}$ (4 elements).
  \item[14.4.] $A = \{\varnothing, \{1\}, \{2\}\}$: $2^A$ has $2^3 = 8$ elements:
    \begin{align*}
      2^A = \bigl\{&\varnothing,\;\{\varnothing\},\;\{\{1\}\},\;\{\{2\}\},\\
      &\{\varnothing,\{1\}\},\;\{\varnothing,\{2\}\},\;\{\{1\},\{2\}\},\;\{\varnothing,\{1\},\{2\}\}\bigr\}
    \end{align*}
  \item[14.5.] $A = \{1, \{3\}, \{1,2\}\}$: $2^A$ has $2^3 = 8$ elements:
    \begin{align*}
      2^A = \bigl\{&\varnothing,\;\{1\},\;\{\{3\}\},\;\{\{1,2\}\},\\
      &\{1,\{3\}\},\;\{1,\{1,2\}\},\;\{\{3\},\{1,2\}\},\;\{1,\{3\},\{1,2\}\}\bigr\}
    \end{align*}
\end{itemize}

%======================================================================
\subsection*{Problem 15}
Prove $\varnothing \subseteq A \cap B \subseteq A \cup B$.

\textit{Proof.} $\varnothing \subseteq S$ for any set $S$ (vacuous truth). If $x \in A \cap B$, then $x \in A$, so $x \in A \cup B$. Hence $A \cap B \subseteq A \cup B$. \qed

%======================================================================
\subsection*{Problem 16}
Describe sets with $U = \mathbb{Z}$, $A = \{2k \mid k \in \mathbb{Z}^+\}$ (positive evens), $B = \{2k-1 \mid k \in \mathbb{Z}^+\}$ (positive odds), $C = \{x \in \mathbb{Z} \mid x < 10\}$.

\begin{itemize}
  \item $\boxed{\overline{A} = \{x \in \mathbb{Z} \mid x \le 0 \text{ or } x \text{ is odd}\}}$
  \item $\boxed{A \triangle B = \mathbb{Z}^+}$ \quad (since $A \cap B = \varnothing$ and $A \cup B = \mathbb{Z}^+$)
  \item $\boxed{\overline{C} = \{x \in \mathbb{Z} \mid x \ge 10\}}$
  \item $\boxed{A \setminus C = \{10, 12, 14, \ldots\}}$ \quad (positive evens $\ge 10$)
  \item $\boxed{C \setminus (A \triangle B) = \{\ldots, -2, -1, 0\}}$ \quad (non-positive integers)
\end{itemize}

%======================================================================
\subsection*{Problem 17}
$A$ = positive evens, $B$ = positive odds, $C$ = positive multiples of 3.

\textbf{(a)}
\begin{itemize}
  \item $A \cap C = \{6, 12, 18, \ldots\}$ = positive multiples of 6.
  \item $B \cup C = \{n \in \mathbb{Z}^+ \mid n \text{ is odd or } 3 \mid n\} = \{1,3,5,6,7,9,11,12,13,\ldots\}$.
  \item $B \setminus C = \{1, 5, 7, 11, 13, 17, \ldots\}$ = positive odd integers not divisible by 3.
\end{itemize}

\textbf{(b)} Yes, the distributive law $A \cap (B \cup C) = (A \cap B) \cup (A \cap C)$ holds for all sets. Since $A \cap B = \varnothing$ (no positive integer is both even and odd), we get $A \cap (B \cup C) = \varnothing \cup (A \cap C) = A \cap C$ = positive multiples of 6. Independently, $A \cap (B \cup C)$ consists of positive even integers that are odd or divisible by 3; the ``odd'' part is impossible, so this is positive evens divisible by 3, i.e., multiples of 6. Both sides agree. $\checkmark$

%======================================================================
\subsection*{Problem 18}
For any set $A$:

$\boxed{A \cap \varnothing = \varnothing}$, \quad $\boxed{A \cup \varnothing = A}$, \quad $\boxed{A \setminus \varnothing = A}$, \quad $\boxed{A \setminus A = \varnothing}$, \quad $\boxed{\varnothing \setminus A = \varnothing}$.

%======================================================================
\subsection*{Problem 19}
Compute:

\begin{itemize}
  \item $\varnothing \setminus \{\varnothing\} = \boxed{\varnothing}$ \quad (nothing to remove from the empty set).
  \item $\{\varnothing\} \setminus \{\varnothing\} = \boxed{\varnothing}$ \quad (the sole element $\varnothing$ is removed).
  \item $\{\varnothing, \{\varnothing\}\} \setminus \varnothing = \boxed{\{\varnothing, \{\varnothing\}\}}$ \quad (nothing is removed).
  \item $\{\varnothing, \{\varnothing\}\} \setminus \{\varnothing\} = \boxed{\{\{\varnothing\}\}}$ \quad (the element $\varnothing$ is removed; $\{\varnothing\}$ remains).
\end{itemize}

%======================================================================
\subsection*{Problem 20}
Prove equivalent characterizations.

\textbf{Part 1:} $A \subseteq B \iff A \cup B = B \iff A \cap B = A$.

\textit{Proof sketch.} $(A \subseteq B) \Rightarrow (A \cup B = B)$: every element of $A$ is already in $B$, so $A \cup B = B$.
$(A \cup B = B) \Rightarrow (A \cap B = A)$: if $x \in A$ then $x \in A \cup B = B$, so $x \in A \cap B$.
$(A \cap B = A) \Rightarrow (A \subseteq B)$: if $x \in A = A \cap B$ then $x \in B$. \qed

\textbf{Part 2:} $A \cap B = \varnothing \iff A \subseteq \overline{B} \iff B \subseteq \overline{A}$.

\textit{Proof sketch.} Disjointness means every element of $A$ is outside $B$ (i.e.\ in $\overline{B}$), and symmetrically every element of $B$ is outside $A$. Cyclic implications close the loop. \qed

\textbf{Part 3:} $A \cup B = U \iff \overline{A} \subseteq B \iff \overline{B} \subseteq A$.

\textit{Proof sketch.} If $A \cup B = U$ and $x \notin A$, then $x \in B$; so $\overline{A} \subseteq B$. If $\overline{A} \subseteq B$ and $x \notin B$, then $x \notin \overline{A}$, i.e.\ $x \in A$; so $\overline{B} \subseteq A$. If $\overline{B} \subseteq A$ then for any $x \in U$, either $x \in B$ or $x \in \overline{B} \subseteq A$, so $A \cup B = U$. \qed

%======================================================================
\subsection*{Problem 21}
Prove $(A \setminus B) \setminus C = A \setminus (B \cup C)$. Is $C \subseteq A$ equivalent?

\textit{Proof.} $x \in (A \setminus B) \setminus C \iff x \in A \land x \notin B \land x \notin C \iff x \in A \land x \notin (B \cup C) \iff x \in A \setminus (B \cup C)$. \qed

The condition $C \subseteq A$ is \textbf{not} equivalent to this identity. The identity $(A \setminus B) \setminus C = A \setminus (B \cup C)$ holds for \emph{all} sets $A, B, C$, regardless of whether $C \subseteq A$.

%======================================================================
\subsection*{Problem 22}
Prove $(A \setminus B) \setminus C = (A \setminus C) \setminus (B \setminus C)$.

\textit{Proof.}
$x \in (A \setminus C) \setminus (B \setminus C) \iff x \in A \land x \notin C \land \lnot(x \in B \land x \notin C)$.
Since $x \notin C$, the condition $\lnot(x \in B \land x \notin C)$ reduces to $x \notin B$.
So this equals $x \in A \land x \notin B \land x \notin C = x \in (A \setminus B) \setminus C$. \qed

%======================================================================
\subsection*{Problem 23}
Prove $1 + 2 + \cdots + n = \frac{n(n+1)}{2}$.

\textit{Proof.} Base: $n=1$: $1 = \frac{1 \cdot 2}{2}$. Inductive step: assume true for $n=k$. Then
$\sum_{i=1}^{k+1} i = \frac{k(k+1)}{2} + (k+1) = \frac{k(k+1)+2(k+1)}{2} = \frac{(k+1)(k+2)}{2}$. \qed

%======================================================================
\subsection*{Problem 24}
Prove $1^2 + 2^2 + \cdots + n^2 = \frac{n(n+1)(2n+1)}{6}$.

\textit{Proof.} Base: $n=1$: $1 = \frac{1 \cdot 2 \cdot 3}{6}$. Inductive step: assume for $k$. Then
$\frac{k(k+1)(2k+1)}{6} + (k+1)^2 = \frac{(k+1)[k(2k+1)+6(k+1)]}{6} = \frac{(k+1)(2k^2+7k+6)}{6} = \frac{(k+1)(k+2)(2k+3)}{6}$. \qed

%======================================================================
\subsection*{Problem 25}
Prove $1 + 4 + 7 + \cdots + (3n-2) = \frac{n(3n-1)}{2}$.

\textit{Proof.} Base: $n=1$: $1 = \frac{1 \cdot 2}{2}$. Inductive step:
$\frac{k(3k-1)}{2} + (3(k+1)-2) = \frac{k(3k-1) + 2(3k+1)}{2} = \frac{3k^2+5k+2}{2} = \frac{(k+1)(3k+2)}{2} = \frac{(k+1)(3(k+1)-1)}{2}$. \qed

%======================================================================
\subsection*{Problem 26}
Prove $\sum_{i=1}^{n} \frac{1}{(2i-1)(2i+1)} = \frac{n}{2n+1}$.

\textit{Proof.} Base: $n=1$: $\frac{1}{1 \cdot 3} = \frac{1}{3}$. Inductive step:
$\frac{k}{2k+1} + \frac{1}{(2k+1)(2k+3)} = \frac{k(2k+3)+1}{(2k+1)(2k+3)} = \frac{2k^2+3k+1}{(2k+1)(2k+3)} = \frac{(k+1)(2k+1)}{(2k+1)(2k+3)} = \frac{k+1}{2k+3}$. \qed

%======================================================================
\subsection*{Problem 27}
Prove $1 + 2 + 2^2 + \cdots + 2^{n-1} = 2^n - 1$.

\textit{Proof.} Base: $n=1$: $1 = 2^1 - 1$. Inductive step:
$(2^k - 1) + 2^k = 2 \cdot 2^k - 1 = 2^{k+1} - 1$. \qed

%======================================================================
\subsection*{Problem 28}
Prove $1 + r + r^2 + \cdots + r^{n-1} = \frac{1-r^n}{1-r}$ for $r \neq 1$.

\textit{Proof.} Base: $n=1$: $1 = \frac{1-r}{1-r}$. Inductive step:
$\frac{1-r^k}{1-r} + r^k = \frac{1-r^k + r^k(1-r)}{1-r} = \frac{1 - r^{k+1}}{1-r}$. \qed

%======================================================================
\subsection*{Problem 29}
Prove $1 + 3 + 5 + \cdots + (2n-1) = n^2$.

\textit{Proof.} Base: $n=1$: $1 = 1^2$. Inductive step:
$k^2 + (2(k+1)-1) = k^2 + 2k + 1 = (k+1)^2$. \qed

%======================================================================
\subsection*{Problem 30}
Prove $2 + 4 + 6 + \cdots + 2n = n^2 + n$.

\textit{Proof.} Base: $n=1$: $2 = 1 + 1$. Inductive step:
$(k^2 + k) + 2(k+1) = k^2 + 3k + 2 = (k+1)^2 + (k+1)$. \qed

%======================================================================
\subsection*{Problem 31}
Prove by induction: $a^{m+n} = a^m \cdot a^n$.

\textit{Proof.} Fix $m$; induct on $n$. Base: $n=0$: $a^{m+0} = a^m = a^m \cdot 1 = a^m \cdot a^0$.
Inductive step: $a^{m+(k+1)} = a^{(m+k)+1} = a^{m+k} \cdot a = (a^m \cdot a^k) \cdot a = a^m \cdot a^{k+1}$. \qed

%======================================================================
\subsection*{Problem 32}
Prove by induction: $(ab)^k = a^k b^k$.

\textit{Proof.} Base: $k=0$: $(ab)^0 = 1 = a^0 b^0$. (Or $k=1$: $(ab)^1 = ab = a^1 b^1$.)
Inductive step: $(ab)^{k+1} = (ab)^k \cdot ab = a^k b^k \cdot ab = a^{k+1} b^{k+1}$ (using commutativity of multiplication). \qed

%======================================================================
\subsection*{Problem 33}
Prove $n^2 > 2n + 1$ for $n \ge 3$.

\textit{Proof.} Base: $n=3$: $9 > 7$. Inductive step: assume $k^2 > 2k+1$ for some $k \ge 3$.
$(k+1)^2 = k^2 + 2k + 1 > (2k+1) + 2k + 1 = 4k + 2 \ge 2(k+1) + 1 + (2k-1) > 2(k+1)+1$
since $2k - 1 \ge 5 > 0$. \qed

%======================================================================
\subsection*{Problem 34}
Prove $2^n > n^2$ for $n \ge 5$.

\textit{Proof.} Base: $n=5$: $32 > 25$. Inductive step: assume $2^k > k^2$ for $k \ge 5$.
$2^{k+1} = 2 \cdot 2^k > 2k^2$. We need $2k^2 \ge (k+1)^2 = k^2 + 2k + 1$, i.e., $k^2 - 2k - 1 \ge 0$, which holds for $k \ge 3$. \qed

%======================================================================
\subsection*{Problem 35}
Prove $a^3 > 3a^2 + 3a + 1$ for $a \ge 4$.

\textit{Proof.} Base: $a=4$: $64 > 48 + 12 + 1 = 61$. $\checkmark$

Inductive step: assume $k^3 > 3k^2 + 3k + 1$ for $k \ge 4$. We show $(k+1)^3 > 3(k+1)^2 + 3(k+1) + 1$.
$(k+1)^3 = k^3 + 3k^2 + 3k + 1 > (3k^2 + 3k + 1) + 3k^2 + 3k + 1 = 6k^2 + 6k + 2$.
We need $6k^2 + 6k + 2 \ge 3k^2 + 9k + 7$, i.e., $3k^2 - 3k - 5 \ge 0$, which holds for $k \ge 2$. \qed

%======================================================================
\subsection*{Problem 36}
Prove $5^{2n} + 3 \cdot 2^{5n-2}$ is divisible by 7 for all positive integers $n$.

\textit{Proof.} Base: $n=1$: $25 + 3 \cdot 8 = 25 + 24 = 49 = 7 \cdot 7$. $\checkmark$

Inductive step: assume $7 \mid (5^{2k} + 3 \cdot 2^{5k-2})$. Then
$5^{2(k+1)} + 3 \cdot 2^{5(k+1)-2} = 25 \cdot 5^{2k} + 32 \cdot 3 \cdot 2^{5k-2}$
$= 25 \cdot 5^{2k} + 25 \cdot 3 \cdot 2^{5k-2} + 7 \cdot 3 \cdot 2^{5k-2}$
$= 25(5^{2k} + 3 \cdot 2^{5k-2}) + 7 \cdot 3 \cdot 2^{5k-2}$.
Both terms are divisible by 7. \qed

%======================================================================
\subsection*{Problem 37}
Prove $3^{n+2} + 4^{2n+1}$ is divisible by 13 for all $n \ge 0$.

\textit{Proof.} Base: $n=0$: $9 + 4 = 13$. $\checkmark$ (Or $n=1$: $27 + 64 = 91 = 7 \cdot 13$.)

Inductive step: assume $13 \mid (3^{k+2} + 4^{2k+1})$. Then
$3^{k+3} + 4^{2k+3} = 3 \cdot 3^{k+2} + 16 \cdot 4^{2k+1}$
$= 3(3^{k+2} + 4^{2k+1}) + 13 \cdot 4^{2k+1}$.
Both terms are divisible by 13. \qed

%======================================================================
\subsection*{Problem 38}
Prove $5^{2n} + 7$ is divisible by 8 for all positive integers $n$.

\textit{Proof.} $5^{2n} = 25^n \equiv 1^n = 1 \pmod{8}$, so $5^{2n} + 7 \equiv 1 + 7 = 8 \equiv 0 \pmod{8}$.

Alternatively by induction: Base: $n=1$: $25 + 7 = 32 = 8 \cdot 4$.
Inductive step: $5^{2(k+1)} + 7 = 25 \cdot 5^{2k} + 7 = 25(5^{2k}+7) - 25 \cdot 7 + 7 = 25(5^{2k}+7) - 168$.
Both $25(5^{2k}+7)$ and $168 = 8 \cdot 21$ are divisible by 8. \qed

%======================================================================
\subsection*{Problem 39}
Prove $3^{3n+1} + 2^{n+1}$ is divisible by 5 for $n \in \mathbb{N}$.

\textit{Proof.} Base: $n=0$: $3 + 2 = 5$. $\checkmark$

Inductive step: $3^{3(k+1)+1} + 2^{k+2} = 27 \cdot 3^{3k+1} + 2 \cdot 2^{k+1}$
$= 27 \cdot 3^{3k+1} + 27 \cdot 2^{k+1} - 25 \cdot 2^{k+1}$
$= 27(3^{3k+1} + 2^{k+1}) - 25 \cdot 2^{k+1}$.
By hypothesis $5 \mid 27(3^{3k+1} + 2^{k+1})$ (since $5$ divides the inner expression) and $5 \mid 25 \cdot 2^{k+1}$. \qed

%======================================================================
\subsection*{Problem 40}
Prove $9^n - 1$ is divisible by 8 for $n \in \mathbb{N}$.

\textit{Proof.} $9^n - 1 = (9-1)(9^{n-1} + 9^{n-2} + \cdots + 1) = 8 \cdot \sum_{i=0}^{n-1} 9^i$.

Alternatively: $9 \equiv 1 \pmod{8}$, so $9^n \equiv 1 \pmod{8}$. \qed

%======================================================================
\subsection*{Problem 41}
Prove $10^n - 1$ is divisible by 9 for $n \in \mathbb{N}$.

\textit{Proof.} $10 \equiv 1 \pmod{9}$, so $10^n \equiv 1 \pmod{9}$, hence $9 \mid (10^n - 1)$.

Alternatively: $10^n - 1 = \underbrace{99\ldots9}_{n} = 9 \cdot \underbrace{11\ldots1}_{n}$. \qed

%======================================================================
\subsection*{Problem 42}
Prove $n^2 - 1$ is divisible by 4 for all odd natural numbers.

\textit{Proof.} Let $n = 2m+1$ for $m \ge 0$. Then $n^2 - 1 = (2m+1)^2 - 1 = 4m^2 + 4m = 4m(m+1)$, which is divisible by 4. \qed

(No induction needed; a direct proof suffices.)

%======================================================================
\subsection*{Problem 43}
Find and prove: $3 + 7 + 11 + \cdots + (4n-1)$.

The sum is $\boxed{n(2n+1)}$.

\textit{Proof.} Base: $n=1$: $3 = 1 \cdot 3$. $\checkmark$

Inductive step: $k(2k+1) + (4(k+1)-1) = 2k^2 + k + 4k + 3 = 2k^2 + 5k + 3 = (k+1)(2k+3) = (k+1)(2(k+1)+1)$. \qed

%======================================================================
\subsection*{Problem 44}
Find and prove: $1 + 5 + 9 + \cdots + (4n-3)$.

The sum is $\boxed{n(2n-1)}$.

\textit{Proof.} Base: $n=1$: $1 = 1 \cdot 1$. $\checkmark$

Inductive step: $k(2k-1) + (4(k+1)-3) = 2k^2 - k + 4k + 1 = 2k^2 + 3k + 1 = (k+1)(2k+1) = (k+1)(2(k+1)-1)$. \qed

%======================================================================
\subsection*{Problem 45}
Find and prove the sum $1\cdot 2 + 3\cdot 4 + 5\cdot 6 + \cdots + (2n-1)\cdot 2n$.

(Interpreting the unclear original as $\sum_{k=1}^{n}(2k-1)(2k)$.)

The sum is $\boxed{\dfrac{n(n+1)(4n-1)}{3}}$.

\textit{Proof.} Base: $n=1$: $1 \cdot 2 = 2$ and $\frac{1 \cdot 2 \cdot 3}{3} = 2$. $\checkmark$

Inductive step: $\frac{k(k+1)(4k-1)}{3} + (2k+1)(2k+2) = \frac{k(k+1)(4k-1) + 6(2k+1)(k+1)}{3}$
$= \frac{(k+1)[k(4k-1)+6(2k+1)]}{3} = \frac{(k+1)(4k^2+11k+6)}{3} = \frac{(k+1)(k+2)(4k+3)}{3}$. \qed

%======================================================================
\subsection*{Problem 46}
Find and prove: $1\cdot 3 + 5\cdot 7 + \cdots + (4n-3)(4n-1)$.

The general term is $(4k-3)(4k-1)$ for $k=1,\ldots,n$. The sum is $\boxed{\dfrac{n(16n^2 - 7)}{3}}$.

\textit{Proof.} Base: $n=1$: $1 \cdot 3 = 3 = \frac{1 \cdot 9}{3}$. $\checkmark$

Inductive step: $\frac{k(16k^2-7)}{3} + (4(k+1)-3)(4(k+1)-1) = \frac{k(16k^2-7)}{3} + (4k+1)(4k+3)$.
Both sides of the target identity $\frac{(k+1)(16(k+1)^2-7)}{3}$ expand to $\frac{16k^3+48k^2+41k+9}{3}$ (routine algebra). \qed

%======================================================================
\subsection*{Problem 47}
Prove $2^n > n^3$ and $n! > 100^n$ for large $n$.

\textbf{(a)} $2^n > n^3$ for $n \ge 10$.

\textit{Proof.} Base: $n=10$: $1024 > 1000$. Inductive step: assume $2^k > k^3$ for $k \ge 10$. Then $2^{k+1} = 2 \cdot 2^k > 2k^3$. We need $2k^3 \ge (k+1)^3$, i.e., $(1+1/k)^3 \le 2$. For $k \ge 10$: $(1+1/k)^3 \le (1.1)^3 = 1.331 < 2$. $\checkmark$ \qed

\textbf{(b)} $n! > 100^n$ for $n \ge 200$ (for instance).

\textit{Proof.} Once $k \ge 200$, we have $k + 1 > 100$, so $(k+1)! = (k+1) \cdot k! > 100 \cdot 100^k = 100^{k+1}$. The base case $200! > 100^{200}$ can be verified using Stirling's approximation. \qed

%======================================================================
\subsection*{Problem 48}
Tower of Hanoi: prove $n$ disks can be moved in $2^n - 1$ moves.

\textit{Proof.} Base: $n=1$: move the single disk directly (1 move $= 2^1 - 1$).
Inductive step: to move $k+1$ disks from peg $A$ to peg $C$, first move the top $k$ disks from $A$ to $B$ ($2^k - 1$ moves by hypothesis), then move the largest disk from $A$ to $C$ (1 move), then move $k$ disks from $B$ to $C$ ($2^k - 1$ moves). Total: $2(2^k-1)+1 = 2^{k+1}-1$. \qed

%======================================================================
\subsection*{Problem 49}
Bernoulli's inequality: $(1+h)^n \ge 1 + nh$ for $h > -1$, $n \in \mathbb{N}$.

\textit{Proof.} Base: $n=1$: $1+h \ge 1+h$. $\checkmark$

Inductive step: $(1+h)^{k+1} = (1+h)^k(1+h) \ge (1+kh)(1+h) = 1 + (k+1)h + kh^2 \ge 1 + (k+1)h$
since $kh^2 \ge 0$. (We used $1+h > 0$ to preserve the inequality direction when multiplying.) \qed

%======================================================================
\subsection*{Problem 50}
Prove Theorem 1 (basic set algebra laws) using membership relation.

``Theorem 1'' typically refers to the fundamental identities:
\begin{enumerate}
  \item Commutativity: $A \cup B = B \cup A$, $A \cap B = B \cap A$.
  \item Associativity: $(A \cup B) \cup C = A \cup (B \cup C)$, similarly for $\cap$.
  \item Distributivity: $A \cup (B \cap C) = (A \cup B) \cap (A \cup C)$, $A \cap (B \cup C) = (A \cap B) \cup (A \cap C)$.
  \item Identity: $A \cup \varnothing = A$, $A \cap U = A$.
  \item Complement: $A \cup \overline{A} = U$, $A \cap \overline{A} = \varnothing$.
  \item Idempotency: $A \cup A = A$, $A \cap A = A$.
  \item De Morgan: $\overline{A \cup B} = \overline{A} \cap \overline{B}$, $\overline{A \cap B} = \overline{A} \cup \overline{B}$.
\end{enumerate}

\textit{Proof sketch.} Each identity is proved by showing $x \in \text{LHS} \iff x \in \text{RHS}$ for arbitrary $x$, translating set operations to logical connectives ($\cup \to \lor$, $\cap \to \land$, $\overline{\cdot} \to \lnot$) and applying standard propositional logic laws. \qed

%======================================================================
\subsection*{Problem 51}
Prove Theorem 2 using membership relation.

``Theorem 2'' typically includes additional identities such as:
absorption ($A \cup (A \cap B) = A$, $A \cap (A \cup B) = A$), involution ($\overline{\overline{A}} = A$), domination ($A \cup U = U$, $A \cap \varnothing = \varnothing$), and $\overline{U} = \varnothing$, $\overline{\varnothing} = U$.

\textit{Proof sketch.} As in Problem 50, each is proved via element-chasing, reducing to tautologies of propositional logic. For example, absorption: $x \in A \cup (A \cap B) \iff x \in A \lor (x \in A \land x \in B) \iff x \in A$. \qed

%======================================================================
\subsection*{Problem 52}
Prove Theorem 2 using Theorem 1 identities as axioms (algebraic proofs, no element-chasing).

\textit{Proof sketch.} For example, absorption: $A \cup (A \cap B) = (A \cap U) \cup (A \cap B) = A \cap (U \cup B) = A \cap U = A$, using identity, distributivity, domination, and identity laws from Theorem 1. Each Theorem 2 identity can be derived purely algebraically from Theorem 1. \qed

%======================================================================
\subsection*{Problem 53}
Using Theorems 1 \& 2, prove identities.

\textbf{(1)} $(A \cap \overline{B} \cap \overline{X}) \cup (A \cap \overline{B} \cap \overline{C} \cap \overline{X} \cap \overline{Y}) \cup (A \cap \overline{X} \cap \overline{A}) = A \cap \overline{B} \cap \overline{X}$.

\textit{Proof sketch.} The third term contains $A \cap \overline{A} = \varnothing$, so it vanishes. The second term is a subset of the first (it has extra constraints $\overline{C}$ and $\overline{Y}$), so $\text{first} \cup \text{second} = \text{first}$, by absorption. Result: $A \cap \overline{B} \cap \overline{X}$. \qed

(The remaining sub-problems follow similarly by applying idempotency, absorption, and complement laws.)

%======================================================================
\subsection*{Problem 54}
Prove/verify set identities. Representative items:

\textbf{(1)} $A \setminus (B \setminus C) = (A \setminus B) \cup C$ --- \textbf{FALSE in general.}

Counterexample: $A = \{1\}$, $B = \varnothing$, $C = \{2\}$. LHS $= \{1\} \setminus \varnothing = \{1\}$. RHS $= \{1\} \cup \{2\} = \{1,2\}$. LHS $\neq$ RHS.

The correct identity is $A \setminus (B \setminus C) = (A \setminus B) \cup (A \cap C)$.

\textbf{(3)} $(A \cup B) \setminus C = (A \setminus C) \cup (B \setminus C)$.

\textit{Proof.} $x \in (A \cup B) \setminus C \iff (x \in A \lor x \in B) \land x \notin C \iff (x \in A \land x \notin C) \lor (x \in B \land x \notin C) \iff x \in (A\setminus C) \cup (B \setminus C)$. \qed

\textbf{(11)} $A \cap (B \mathbin{\triangle} C) = (A \cap B) \mathbin{\triangle} (A \cap C)$.

\textit{Proof.} Both sides reduce to $(x \in A \land x \in B \land x \notin C) \lor (x \in A \land x \notin B \land x \in C)$. \qed

(Other sub-problems in 54 follow the same element-chasing technique, translating $\setminus$ as $\cap\;\overline{(\cdot)}$ and using propositional logic equivalences.)

%======================================================================
\subsection*{Problem 55}
Show $A \cup (B \setminus C) \neq (A \cup B) \setminus (A \cup C)$ in general.

Counterexample: $A = \{1\}$, $B = \{1,2\}$, $C = \{1\}$.
LHS: $B \setminus C = \{2\}$, so $A \cup (B \setminus C) = \{1,2\}$.
RHS: $A \cup B = \{1,2\}$, $A \cup C = \{1\}$, so $(A \cup B) \setminus (A \cup C) = \{2\}$.
$\{1,2\} \neq \{2\}$. \qed

%======================================================================
\subsection*{Problem 56}
Prove $A \mathbin{\triangle} X = B$ has a unique solution.

The unique solution is $\boxed{X = A \mathbin{\triangle} B}$.

\textit{Proof.}
\textbf{Existence:} Let $X = A \mathbin{\triangle} B$. Then $A \mathbin{\triangle} X = A \mathbin{\triangle} (A \mathbin{\triangle} B) = (A \mathbin{\triangle} A) \mathbin{\triangle} B = \varnothing \mathbin{\triangle} B = B$. $\checkmark$

\textbf{Uniqueness:} If $A \mathbin{\triangle} X = B$, apply $A \mathbin{\triangle}$ to both sides: $A \mathbin{\triangle} (A \mathbin{\triangle} X) = A \mathbin{\triangle} B$, giving $X = A \mathbin{\triangle} B$. (We used associativity of $\triangle$ and $A \mathbin{\triangle} A = \varnothing$, $\varnothing \mathbin{\triangle} X = X$.) \qed

%======================================================================
\subsection*{Problem 57}
Do there exist sets $A, B$ such that $|A \times B|$ equals:

\textbf{57.1.} 12 elements? \textbf{YES.} For example $A = \{1,2,3\}$, $B = \{a,b,c,d\}$; $|A| \cdot |B| = 3 \cdot 4 = 12$.

\textbf{57.2.} 13 elements? \textbf{YES.} Since $13$ is prime, $|A| \cdot |B| = 13$ forces $|A|=1, |B|=13$ or vice versa. For example $A=\{1\}$, $B$ any 13-element set.

$\boxed{\text{57.1: Yes (e.g., } 3 \times 4 = 12\text{). \quad 57.2: Yes (e.g., } 1 \times 13 = 13\text{).}}$

%======================================================================
\subsection*{Problem 58}
Prove $A \times B = B \times A \implies A = B$ (assuming both non-empty).

\textit{Proof.} Assume $A, B \neq \varnothing$ and $A \times B = B \times A$.
To show $A \subseteq B$: let $a \in A$, pick any $b \in B$. Then $(a,b) \in A \times B = B \times A$, so $a \in B$.
To show $B \subseteq A$: let $b \in B$, pick any $a \in A$. Then $(b,a) \in B \times A = A \times B$, so $b \in A$.
Hence $A = B$. (If either set is empty, both products are $\varnothing$ regardless of $A = B$.) \qed

%======================================================================
\subsection*{Problem 59}
List elements of $A \times B$ for $A = \{a,b\}$, $B = \{1,2,3,4\}$.

\[
  \boxed{A \times B = \{(a,1),(a,2),(a,3),(a,4),(b,1),(b,2),(b,3),(b,4)\}}
\]
(8 elements.)

%======================================================================
\subsection*{Problem 60}
Prove Cartesian product distributes over set operations.

\textbf{(1)} $A \times (B \cup C) = (A \times B) \cup (A \times C)$.

\textit{Proof.} $(x,y) \in A \times (B \cup C) \iff x \in A \land (y \in B \lor y \in C) \iff (x \in A \land y \in B) \lor (x \in A \land y \in C) \iff (x,y) \in (A\times B) \cup (A \times C)$. \qed

\textbf{(2)} $(B \cup C) \times A = (B \times A) \cup (C \times A)$. Same proof on the first component.

\textbf{(3)} $A \times (B \cap C) = (A \times B) \cap (A \times C)$.

\textit{Proof.} $(x,y) \in \text{LHS} \iff x \in A \land y \in B \land y \in C \iff (x \in A \land y \in B) \land (x \in A \land y \in C)$. \qed

\textbf{(4)} $(B \cap C) \times A = (B \times A) \cap (C \times A)$. Mirror of (3).

\textbf{(5)} $A \times (B \setminus C) = (A \times B) \setminus (A \times C)$.

\textit{Proof.} $(x,y) \in \text{LHS} \iff x \in A \land y \in B \land y \notin C$. $(x,y) \in \text{RHS} \iff (x \in A \land y \in B) \land \lnot(x \in A \land y \in C)$. Given $x \in A$, the second condition reduces to $y \notin C$. \qed

\textbf{(6)} $(B \setminus C) \times A = (B \times A) \setminus (C \times A)$. Mirror of (5).

%======================================================================
\subsection*{Problem 61}
Prove $(A \cap B) \times (C \cap D) = (A \times C) \cap (B \times D)$.

\textit{Proof.}
$(x,y) \in (A \cap B) \times (C \cap D)$
$\iff x \in A \land x \in B \land y \in C \land y \in D$
$\iff (x \in A \land y \in C) \land (x \in B \land y \in D)$
$\iff (x,y) \in (A \times C) \cap (B \times D)$. \qed

%======================================================================
\subsection*{Problem 62}
Prove $A \subseteq B \land C \subseteq D \implies A \times C \subseteq B \times D$.

\textit{Proof.} Let $(x,y) \in A \times C$. Then $x \in A \subseteq B$ gives $x \in B$, and $y \in C \subseteq D$ gives $y \in D$. So $(x,y) \in B \times D$. \qed
